% a01_H0rel part 6b: the steps of the upper bound
\paragraph{The steps of the upper bound.}
\begin{description}[leftmargin=1.5em,style=nextline]
\item[{D1 (large/small split; \cite[(2.2)--(2.3)]{Ba2})}] At each step the fields $B,A,\psi,\phi$ are tested; each test is tied to a positive term of the action and a violation yields $e^{-c_0^{\mathrm{Ba}}p(\ell_i)^2}$; with the Dirichlet condition the matter Gaussian only gains the positive face terms, which lower the integrand, and every printed large-field factor is monotone in them.
\item[D2 (the face family)] At every scale $0\le i\le K_f$ the face test with threshold $\sigma_i$ is added to the six families (\cref{def:a01:facefam}); a violation is paid $e^{-p_i^2/4}$ by \cref{lem:a01:F541}(a) from the $\theta_2$-share of the scale-$i$ form, the $(1-\theta_2)$-share remaining for D12. The family is forced: without it the background of face data of size $M_i$ has a face gradient of order $M_i$ (\cref{prop:a01:legsd}, remark), which violates the fine-field class \cite[(2.55)]{Ba2}.
\item[{D3 (region geometry; \cite[(2.4)--(2.10), (2.43)]{Ba2})}] As said above, verbatim with seven families. The pair $B^i(y)\cup B^i(y')$ of \cref{lem:a01:FPC}'s payment is one block here, so no clean-pair condition arises; tested blocks are clean (no Neumann cut through $B^i(y')$) by the nesting rules.
\item[{D4 (scalar large fields; \cite[(2.11)--(2.13)]{Ba2})}] $\lambda(\eps)|\phi|^4+\frac12\dm\eps^2|\phi|^2\ge-O(\eps^{\beta_0})$ and $\ge\frac12\lambda(\eps)|\phi|^4$ for $|\phi|^2\ge O(1+\ln\eps^{-1})$: unchanged, $\dm$ being King's free-boundary counterterm, a vertex, weight-independent.
\item[{D5 (translation and splitting; \cite[(2.18)--(2.19)]{Ba2})}] The photon translation with the Dirichlet-conditioned $C^{(0)}_{\Lambda_0}$ is \cref{lem:a01:PE}(c); the scalar analogue uses the operators on sub-regions with \cref{lem:a01:A1,lem:a01:A23,lem:a01:A4,lem:a01:DV,lem:a01:DVloc} at every point and \cref{lem:a01:A5,lem:a01:legs} in the restricted form (depth $\ge R_{\mathrm{Ba}}$, which is where Ba{\l}aban's splitting evaluates them: his vertex sets are $\Lambda_7$-level and his propagator domains $\Lambda_2$-level), \cref{thm:a01:B2} for the perturbative terms on these sub-domains and \cref{prop:a01:dGD}(i), (ii), (iv), (vii) for the changes of domain and of boundary condition. No face hypothesis on the splitting domains is needed on the cube.
\item[{D6 (expansion to order $\bar n$, cumulants, removal of the interaction from $\Lambda_7^c$; \cite[(2.22)--(2.28), Prop.~2.6]{Ba2})}] \Cref{prop:a01:TD}(c) and \cref{thm:a01:B2} with the difference lines of the pair of propagator domains of \cite[Prop.~2.6]{Ba2} as external pairs (the telescoped families are finite signed sums of members with $s\ge1$ difference objects, each admissible with the gain distance $\dist(\cdot,\Xl\setminus\Omega_2)$, so the extra factor is $e^{-5\delta_cr(\ell_k)}=O(\ell_k^{\kappa})$ at the margin $>5r(\ell_k)$ of Ba{\l}aban's Proposition~2.6: this is the interaction-term bound \cref{hyp:k01:BALint} on the cube, with the printed $e^{-5\delta_0r(\ell_k)}$ replaced by $e^{-5\delta_cr(\ell_k)}$); the replacements \cite[(2.108)--(2.115)]{Ba2} by \cref{prop:a01:dGD}(iii), (v), (vi).
\item[{D7 (the collar quartic; \cite[Prop.~2.1, (2.57), (2.79)]{Ba2})}] \Cref{thm:a01:B2} on the propagator domain of \cite[(2.56)]{Ba2} keeps $-\lambda(\ell_i)\sum_{x\in B^i(\mathrm{collar})}\vth^d|\phi^{(i)}(x)|^4$ exactly. The step \cite[(2.79)]{Ba2} rests on \cite[(2.65)]{Ba2}, i.e.\ on $|\phi^{(i)}(x)|\ge\frac12|\phi_i(y)|$ for $x\in B^i(y)$, $y\in R_s^{(i)}$. (b) \emph{The near layers:} for $y\in R_s^{(i)}$ within $\ell^\infty$-distance $R_0$ of the face layer, \cref{lem:a01:F541}(e) supplies a witness (a gradient violation or a face-family element) within distance $dR_0+1$, paid by its own family. (c) \emph{All other blocks:} for $y$ at distance $>R_0$ from $\DD$, the background on $B^i(y)$ differs from the bulk background by the image effect, bounded by $C'e^{-\delta_1(R_0-2)}M_i\le\frac14|\phi_i(y)|$ on the box $\Omega_b\defeq\Xl\cap(\text{the }\ell^\infty\text{ block ball of radius }\lfloor\dep(x)/2\rfloor\text{ about }y)$ (\cref{lem:a01:FL}, Step~1's box-versus-free factor at $A=0$ and Step~3 at regular $A$) plus the box-versus-domain difference (a cut-off and commutator argument at the Neumann faces of $\Omega_b$, as in \cref{thm:a01:B2}(1), of size $CM_ie^{-(\delta_2/2)(\dep(x)/2-8)}\le p_i$ at the D7 points, which lie several $r$-layers inside the propagator domain), with $R_0$ the smallest integer $\ge3$ with $C'e^{-\delta_1(R_0-2)}\le\frac14$; $R_0$ depends on $(d,L,a,\Wt,c_V)$ only and is fixed before Ba{\l}aban's $M$ ($s_b=LM\ge dR_0+L+1$). Under the caps $(C_B+3)\lambda(\ell_i)^{1/4}\le\frac18$, $C'(V_i+\theta_i)\le\frac18$, this gives $|\phi^{(i)}(x)|\ge\frac12|\phi_i(y)|$, hence \cite[(2.79)]{Ba2} with $\frac12\to\frac1{16}$, and the exponent of $R_s$ in \cite[(3.41)]{Ba2} changes only in its constant.
\item[{D8 (the backgrounds in the regular and fine-field classes; \cite[Lemmas~2.3--2.4, Prop.~2.2]{Ba2})}] $A^{(i)}$: \cref{lem:a01:PE}(a), (d) and \cref{lem:a01:REG}. $\phi^{(i)}$: \cite[Prop.~2.2]{Ba2} is \cref{lem:a01:A5} (every bond, face bonds included); the fine-field class \cite[(2.55)]{Ba2}, $|D_{\tilde A}\phi^{(i)}(b)|\le c_1'p_{i-1}$ on $B^i(\Lambda_7^{(i)})$, is \cref{prop:a01:legsd} with the cap $|\phi_i(y_0)|\le\sigma_i+\dist(y_0,\Lone^{(i)})p_i$ for $\dist\le r_i$ (transport along a block path to the face layer inside $\Lambda_7^{(i-1)\prime}\setminus S_i$, since $\dist(\Lambda_7^{(i)},X_i)>7r_i$) and $e^{-\delta_1r_i}M_i\le p_i$ for $\dist>r_i$; the product $(\sigma_i+tp_i)e^{-\delta_1t}\le Cp_i$. At the collar cubes $\square\subset B^i(\Lambda_7^{(i-1)\prime})$ of \cite[Prop.~2.1]{Ba2}, where the data of \cite[(2.56)]{Ba2} are cut one previous-scale layer away and the scale-$i$ tests are known only off the bad set, the cap is inherited from step $i-1$ (\cref{lem:a01:IC}), inserted as a characteristic function of $\phi_i$ as \cite[(2.55)]{Ba2} is carried to step $i$; with it, for every fine bond $b$ within $1$ of $\square$: $|D_A\phi^{(i)}_\Omega(b)|\le C_Gp_i$ and $|\phi_i(y)-(Q\phi^{(i)}_\Omega)(y)|\le C_Kp_i$ at the block $y$ of $\square$, by \cref{prop:a01:legsd} on the King box of radius $\lfloor m_i/2\rfloor$ about $\square$ inside the propagator domain ($m_i\defeq r(\ell_{i-1})/L\ge70$), the truncation entering as the tail $Cc_1M_ie^{-\delta_cr(\ell_{i-1})/L}\le p_i$ and the box-versus-domain difference as in D7(c); the block clause from \cref{lem:a01:EB}. Every use of \cite[(2.55)]{Ba2} in \cite{Ba2} is polynomial in $p$ and paid by factors $e^{-cp^2}$.
\item[{D9 (fluctuation fields small; \cite[Lemmas~2.5, 2.7]{Ba2})}] $|A'|\le O(1)p$ \cite[(2.86)]{Ba2} by \cref{lem:a01:PE}; $|\phi'|\le O(1)p$ \cite[(2.114)]{Ba2} from the block tests and the covariance bounds (\cref{lem:a01:A23}(c), whose lower bound at the weight is \cref{lem:a01:329} plus the coarse block Poincar\'e inequality); near the face \cref{lem:a01:F541}(b)--(d) with the corrected threshold $r(\ell_i)$ (the contour of \cite[(5.44)]{K2} must stay where the gradient test holds, which a path of length $\le r_i$ from $x\in\Lambda_4^{(i)}$ does, since $\dist(\Lambda_4^{(i)},\Lambda_0^{(i)c})\ge4r_i$). \Cref{lem:a01:F541}(d) supplies \cite[(2.113)]{Ba2} at distance $>r_i$ from the face.
\item[{D10 (the $Z$-factors and the tail removal; \cite[(2.92)--(2.109)]{Ba2})}] $\det(C^{(i)})^{\pm1/2}$ with $C^{(i)}$ bounded (\cref{lem:a01:A23}); the removal of the tail field $\tilde B^{(i)}$ of \cite[(2.89)--(2.90)]{Ba2} from the $Z$-factors, the lines and legs and the quadratic form costs $C\ell_i^{\kappa'}|\Lambda^{(i)}|$: the matter form by the envelope argument at the regular backgrounds, the determinant by the exact block split $\ln\det\Mop=\ln\det(-\Delta_A+V)+\ln\det(1+a_iQ(-\Delta_A+V)^{-1}Q^*)$ with the coarse determinant bounded in summed form and the block-free one by its closed-loop representation at the two regular backgrounds, and the perturbative part by \cref{thm:a01:B2} with the tail as an external photon field; on sub-domain histories the tail lies at distance $>r(\ell_i)/2$ from every place where the domain operators differ from the cube operators, so the domain remainders are $O(\ell_i^\kappa)e^{-cr(\ell_i)}$ times the volume.
\item[{D11 (the last step; \cite[(3.11)--(3.12)]{Ba2})}] $G'_k\ge\gamma_1\lmu^2\ell_k^2$ (photon, \cref{lem:a01:PE}(e)); $G''_k\ge\gamma_1m^2\ell_k^2$ on $\Lambda_5^{(k)}$ from $V\ge m^2$ and the $\ell^2$ bounds of \cref{lem:a01:A23}, the Dirichlet terms only adding; $\norm{H'_k},\norm{H''_k}\le O(1)e^{-\delta_1r(\ell_k)}$ requires only $\ell^2$ bounds (\cref{lem:a01:A23}(a)).
\item[{D12 (positivity; \cite[Prop.~3.1, (3.26)]{Ba2})}] For every configuration in the support of the characteristic functions of \cite[(3.21)]{Ba2} (the fine-field class on each collar), in the multiscale form \cite[(3.23)--(3.25)]{Ba2}: \cite[(3.26)]{Ba2} on the $(1-\theta_2)$-share follows from the Neumann separation \cite[(3.27)--(3.28)]{Ba2} (printed algebra, weight-free) and \cref{lem:a01:329} (the per-scale lower bound \cite[(3.29)]{Ba2} at the weight, whose $\phi$-independent error $O(\ell_k^{\kappa_0})|\Lambda_k|$ is carried on that support exactly as in print), with the Dirichlet faces only adding; the photon exponents are unchanged (\cref{lem:a01:PE}(e)); the face family is paid on the $\theta_2$-share by the composite clause of \cref{lem:a01:FPC}, on blocks of the index set, and payments of different scales are on disjoint fine sets and add. The convergence factors \cite[(3.31)--(3.34)]{Ba2} follow with $(1-\theta_2)$ on the two matter positivity exponents, the photon exponents unchanged, plus $\exp(-\frac14p_i^2|\text{face violations at scale }i|)$ at every scale.
\item[{D13 (integration of the block fields; \cite[(3.32)--(3.38)]{Ba2})}] The normalisation identities are the minimiser identity; the exceptional integrals over $|\phi_{i+1}(y)-(Q\phi_i)(y)|>p$ as printed; after the rescaling, \cite[(3.37)]{Ba2} is $\int d\phi_0\exp(-\frac12\ip{\phi_0}{(-\Delta^D_{\tilde A}+V)\phi_0})$ with the Dirichlet Laplacian and $V\ge m^2>0$, and the diamagnetic inequality of \cite[Sect.~3.1, (3.4)--(3.7)]{BFS1} in its site-diagonal form (the resolvent Neumann series with the site-diagonal $\vth d_x+\vth^3V(x)+$ the face terms, every walk weight non-negative and $\norm{U(\tilde A)}=1$, so $|(-\Delta^D_{\tilde A}+V)^{-1}(x,y)|\le(-\Delta^D+V)^{-1}(x,y)$, and the interpolation of \cite[Sect.~3.3, (3.16)--(3.17)]{BFS1}) gives $\le\int\exp(-\frac12\ip\phi{(-\Delta^D+V)\phi})=e^{E_{0,s}}$.
\item[{D14 (combinatorics; \cite[Sect.~3C, (3.51)--(3.54)]{Ba2})}] Purely combinatorial: $|\Lambda_0^{(i)c}|\le O(1)r(\ell_i)^d(|\mathcal C_0|+\dots+|\mathcal C_i|)$ and $-\sum_i(c_0'p_i^2-O(1)r(\ell_i)^{d+1})|\mathcal C_i|\le0$ if $2p>(d+1)\varrho$ (here $12>8$) and $b_0$ is large, with the product partition into boxes of large blocks in place of \cite[(3.43)--(3.45)]{Ba2} and the cube count $|\mathcal C_i|\le5^d|F_i^{\mathrm{paid}}|$ of \cref{lem:a01:PP5}, the paid part of the family comprising Ba{\l}aban's six families ($R_s$ restricted to its far elements) and the face family (the near $R_s$ elements are charged to their witnesses of D7(b) and not counted); the constant becomes $c_0'\defeq5^{-d}\min\{\frac18(1-\theta_2)a,\frac14(1-\theta_2)\gamma_0^w,\frac14,\frac1{16}\}$, the region entropy $2^{7|\Lambda_0^{(i)c}|}$ is absorbed in the $O(1)(1+\ln\ell_i^{-1})|\Lambda_0^{(i)c}|$ of \cite[(3.47)]{Ba2}. We use half of each payment for \cite[(3.54)]{Ba2} and keep $e^{-\frac12c_0'p_i^2|\mathcal C_i|}$.
\item[{D15 (final step; \cite[(2.116)--(2.118)]{Ba2})}] $P^{(K_f)}\le O(\ell_f^{\beta_0})|\Lambda_7^{(K_f)}|$, $\sum_{i\le K_f}O(\ell_i^{\kappa'})|X^{(i)}|\le O(1)|\CJ|$ (geometric, $\kappa'>d$), and the final integral with the region entropy gives $\exp(O(1)|X^{(K_f)}|)$, $|X^{(K_f)}|=\ell_f^{-d}|\CJ|$: \cref{prop:a01:TD}(c).
\end{description}

\begin{proposition}[upper line]\label{prop:a01:upper}
Under the hypotheses of \cref{prop:a01:lower}, $\Zk(s)\le e^{a}\Ik j(s)+e^{-\frac12c_0'p(\eta)^2+E_+|\CJ|}$, with $a$ as there, $c_0'$ the constant of D14 and $E_+$ depending on $\eps_0^{\mathrm{Ba}}$ and King's data only.
\end{proposition}
\begin{proof}
In the inductive formula \cite[(2.43)]{Ba2} at scale $j$ (seven families at every scale $i\le j$, the gradient family at scale $j$ imposed on every block bond of the cube), the term with no large-field block at any scale $i\le j$ is $\int\chi_j\lambda_j\,e^{-S^{(j)}_k[A_\theta]+\sum_{i<j}O(1)\ell_i^{\kappa'}|X^{(i)}|}$ (D1--D10), with Ba{\l}aban's cut background $A_\theta=a_j\ell_j^{-2}\zeta^{(j)}_\theta G_jQ^*A_j$ in every scalar expression \cite[(2.44), (2.50)]{Ba2}; the difference $|S^{(j)}_k[A^{(j)}]-S^{(j)}_k[A_\theta]|\le C\eta^{s'}|X^{(j)}|$ on $\supp\chi_j$ (the matter form by the envelope argument between two regular backgrounds; the determinant by the block split of D10 at the two backgrounds; the perturbative part by \cref{prop:a01:TD} with the background-type photon field $A^{(j)}-A_\theta$), so the term is $\le e^{a}\Ik j(s)$ after enlarging $C$ in $a$. Every other term has a large-field block at some scale $i\le j$ and carries the kept factor $e^{-\frac12c_0'p_i^2}\le e^{-\frac12c_0'p(\eta)^2}$ ($p$ is decreasing in the scale). Continue the RG for these terms to the final scale $\ell_f\asymp\eps_0^{\mathrm{Ba}}$ (D11--D15, with the face family imposed and paid at each step, D2); the sum over all admissible histories with the remaining half of the payments and \cite[(3.54)]{Ba2} is $\le e^{O(1)|X^{(K_f)}|}=e^{E_+|\CJ|}$, $E_+=O(\ell_f^{-d})$ (\cite[(2.117), (3.42)--(3.47)]{Ba2} with (3.54)). The constants are those of D1--D15: independent of $k$, $j$ and $\vth$, each step's constants being so (the payments $e^{-p_i^2/4}$ are exact for every scale by the choice of $\sigma_i$; at every scale the lattice regime is the only regime on the cube, and \cref{lem:a01:DV,lem:a01:A5} cover it uniformly in $\vth\le1$).
\end{proof}

